\chapter{Extraits du code du prototype LSM}

Cette annexe rassemble des extraits représentatifs du code réellement utilisé
dans le prototype LSM.

\section{Extraction des points de pose et des mains}

Le traitement vidéo conserve les dernières coordonnées disponibles lorsqu'un
repère n'est temporairement pas détecté. Les 33 points du corps et les 21 points
de chaque main forment 225 valeurs par image.

\begin{verbatim}
if results.pose_landmarks:
    corps = np.array([
        [lm.x, lm.y, lm.z]
        for lm in results.pose_landmarks.landmark
    ])
    derniere_pose = corps
else:
    corps = derniere_pose
    n_sans_detection += 1

vecteur = np.concatenate([
    corps.flatten(),
    main_g.flatten(),
    main_d.flatten()
])
frames.append(vecteur)
\end{verbatim}

\section{Rééchantillonnage et normalisation}

Les séquences sont ramenées à 30 frames. À chaque frame, les coordonnées sont
centrées sur le milieu des épaules et divisées par la distance entre celles-ci.

\begin{verbatim}
indices = np.linspace(0, len(sequence) - 1, n_cible)
sequence = np.array([
    sequence[int(round(index))] for index in indices
])

for frame_index in range(sequence.shape[0]):
    epaule_g = sequence[frame_index, 33:36]
    epaule_d = sequence[frame_index, 36:39]
    centre = (epaule_g + epaule_d) / 2.0
    distance = np.linalg.norm(epaule_g - epaule_d)
    if distance < 1e-6:
        distance = 1.0
    for point_index in range(0, 225, 3):
        sequence[frame_index, point_index:point_index + 3] = (
            sequence[frame_index, point_index:point_index + 3] - centre
        ) / distance
\end{verbatim}

\section{Conversion du modèle Keras en TFLite}

Le script de conversion conserve les poids appris, active l'optimisation
TFLite et vérifie l'artefact avec un interpréteur avant son écriture.

\begin{verbatim}
model = tf.keras.models.load_model(args.model, compile=False)
model = make_tflite_compatible_model(model, tf)
converter = tf.lite.TFLiteConverter.from_keras_model(model)
converter.optimizations = [tf.lite.Optimize.DEFAULT]
model_content = converter.convert()

interpreter = tf.lite.Interpreter(model_content=model_content)
interpreter.allocate_tensors()
args.output.parent.mkdir(parents=True, exist_ok=True)
args.output.write_bytes(model_content)
\end{verbatim}

\section{Inférence TFLite}

L'interpréteur est initialisé une fois au démarrage du backend. Son accès est
sérialisé pendant l'inférence et la classe présentant la probabilité maximale
est renvoyée avec son score.

\begin{verbatim}
with _INTERPRETER_LOCK:
    INTERPRETER.set_tensor(INPUT_DETAILS["index"], batch)
    INTERPRETER.invoke()
    probabilites = INTERPRETER.get_tensor(
        OUTPUT_DETAILS["index"]
    )[0]

index = int(np.argmax(probabilites))
return {
    "signe": CLASSES[index],
    "confiance": float(probabilites[index]),
}
\end{verbatim}

\section{Endpoint FastAPI}

L'endpoint reçoit le clip vidéo sous forme de fichier multipart et délègue le
traitement au module d'inférence. Les vidéos non exploitables sont renvoyées
comme erreurs de requête.

\begin{verbatim}
@app.post("/predire")
def predire_endpoint(file: UploadFile = File(...)):
    video_bytes = file.file.read()
    try:
        return predire(video_bytes)
    except InvalidVideoError as error:
        raise HTTPException(
            status_code=422, detail=str(error)
        ) from error
\end{verbatim}

\section{Capture webcam et appel de l'API}

Le navigateur enregistre un clip à partir du flux webcam puis l'envoie à
l'endpoint de prédiction avec \texttt{FormData}.

\begin{verbatim}
const chunks = [];
recorder.addEventListener("dataavailable", (event) => {
  if (event.data.size > 0) chunks.push(event.data);
});
recorder.start();
window.setTimeout(() => {
  if (recorder.state === "recording") recorder.stop();
}, captureDurationMs());

const formData = new FormData();
formData.append("file", clip, "capture.webm");
const response = await fetch("/predire", {
  method: "POST",
  body: formData,
});
\end{verbatim}

\section{Synthèse vocale et création du fichier WAV}

Le texte prédit est vocalisé dans le navigateur avec Transformers.js et le
modèle ONNX malgache. Le signal produit est converti en WAV PCM pour le lecteur
audio HTML.

\begin{verbatim}
const { pipeline, env } = await import(
  "https://cdn.jsdelivr.net/npm/@xenova/transformers"
);
env.backends.onnx.wasm.numThreads = 1;
const synthesizer = await pipeline(
  "text-to-speech",
  "elloza/mms-tts-mlg-onnx",
  { quantized: true }
);
const speech = await synthesizer(text);
audioResult.src = URL.createObjectURL(
  makeWav(speech.audio, speech.sampling_rate)
);
\end{verbatim}

\endinput
